% Generated draft: requires mathematical and source review.
\documentclass[11pt]{article}
\usepackage[margin=25mm]{geometry}
\usepackage{fontspec}
\setmainfont{Latin Modern Roman}
\usepackage{amsmath,amssymb,mathtools,unicode-math}
\setmathfont{Latin Modern Math}
\usepackage{enumitem,needspace,xurl,hyperref}
\setlength{\parindent}{0pt}
\setlength{\parskip}{5pt}
\setlength{\emergencystretch}{3em}
\newcommand{\N}{\mathbb{N}}
\newcommand{\Z}{\mathbb{Z}}
\newcommand{\Q}{\mathbb{Q}}
\newcommand{\R}{\mathbb{R}}
\newcommand{\C}{\mathbb{C}}
\newcommand{\F}{\mathbb{F}}
\newcommand{\A}{\mathbb{A}}
\newcommand{\PP}{\mathbb{P}}
\newcommand{\E}{\mathbb{E}}
\newcommand{\eps}{\varepsilon}
\newcommand{\dd}{\,\mathrm{d}}
\begin{document}
\section*{UnsolvedMath ID: 30003788}
\textbf{Source problem code:} catalogue entry
\subsection{Definitions and mathematical statement}
The supplied record poses four questions. Several parameters and technical conventions are missing; the formulations below expose these omissions rather than resolve them by assumption.

\begin{enumerate}
\item \textbf{Coarse embedding of function spaces.} For a specified exponent \(1\leq p<\infty\) and scalar field \(\mathbb K\in\{\R,\C\}\), let \(L_p([0,1];\mathbb K)\) be the space of Lebesgue-measurable functions modulo equality almost everywhere with norm
\[
\|f\|_p=\left(\int_0^1 |f(t)|^p\,\dd t\right)^{1/p},
\]
and let \(\ell_p(\mathbb K)\) be the space of sequences \(a=(a_j)_{j\geq1}\) with \(\|a\|_p=(\sum_{j\geq1}|a_j|^p)^{1/p}<\infty\). A map \(F:X\to Y\) between metric spaces is a coarse embedding if there are nondecreasing functions \(\rho_-,\rho_+:[0,\infty)\to[0,\infty)\), with \(\rho_-(r)\to\infty\) as \(r\to\infty\), such that
\[
\rho_-(d_X(x,x'))\leq d_Y(F(x),F(x'))\leq\rho_+(d_X(x,x'))
\quad(x,x'\in X).
\]
The question is whether such a map
\[
F:L_p([0,1])\longrightarrow\ell_p
\]
exists. The source does not specify the range or quantifier for \(p\), the scalar field, or whether endpoint or quasi-Banach cases are intended. The displayed definitions cover the stated conventional finite Banach-exponent setting only; they do not supply a missing restriction on the source question.

\item \textbf{Warped cones and walls.} The source asks whether there exist a free group \(\Gamma\), a manifold \(M\), a measure \(\mu\) on \(M\), and a free measure-preserving action \(\Gamma\curvearrowright(M,\mu)\) whose warped cone admits a wall structure. Measure preservation means \(\mu(gA)=\mu(A)\) for every measurable \(A\subseteq M\) and every \(g\in\Gamma\). Pointwise freeness means that \(gx=x\) implies \(g=e\); essential freeness would require this only for almost every \(x\).

For orientation, given a metric \(d_M\), a finite symmetric generating set \(S\) of \(\Gamma\), and suitable action regularity, the usual warped metric on a level \(t\geq1\) is the largest metric \(d_t\) satisfying
\[
d_t(x,y)\leq t\,d_M(x,y),\qquad d_t(x,sx)\leq1
\quad(x,y\in M,\ s\in S).
\]
A wall on a set \(X\) is an unordered partition \(\{H,X\setminus H\}\) into two nonempty parts. In a measured wall structure, a measure on an appropriate measurable collection of walls gives a wall pseudometric by measuring the walls that separate two points. Merely possessing walls does not state any compatibility with the warped metric.

The source leaves unspecified the rank of \(\Gamma\), pointwise versus essential freeness, the dimension and compactness of \(M\), the metric, the measure, action regularity, the full-cone versus level-family convention, and the required meaning and metric compatibility of a wall structure. Thus these definitions explain the question but do not determine a unique precise existence assertion.

\item \textbf{Slow growth of a box space.} Under the usual normal-filtration convention, take a finitely generated residually finite group \(\Gamma\), a finite symmetric generating set \(S\), and finite-index normal subgroups
\[
\Gamma\geq N_1\geq N_2\geq\cdots,
\qquad \bigcap_{m\geq1}N_m=\{e\}.
\]
Its box space is the coarse disjoint union of the finite Cayley graphs \(\operatorname{Cay}(\Gamma/N_m,\overline S)\), each with its graph metric, with distinct components separated by distances tending to infinity. Bounded successive indices would mean
\[
\exists C<\infty\ \forall m\geq1:\quad [N_m:N_{m+1}]\leq C.
\]
The source asks for a slowly growing box space with bounded successive indices. It does not define ``slowly growing,'' identify the quantity whose growth is controlled, or prescribe a comparison rate. It also does not specify the underlying group or whether additional filtration conditions are required. Consequently, the bounded-index condition can be written explicitly, but the requested growth condition cannot be supplied from this record.

\item \textbf{Superexpanders from special linear groups.} Write
\[
G_m=\operatorname{SL}(n_m,\mathbb F_{p_m})
=\{A\in M_{n_m}(\mathbb F_{p_m}):\det A=1\},
\]
where \(p_m\) is prime and \(n_m\) is a positive integer. For a symmetric generating set \(S_m\subseteq G_m\setminus\{e\}\), the undirected Cayley graph has vertex set \(G_m\) and edges \(\{g,gs\}\) for \(g\in G_m\), \(s\in S_m\). The requested fixed-cardinality condition is
\[
\exists k<\infty\ \forall m:\quad |S_m|=k,
\qquad \langle S_m\rangle=G_m.
\]
The question asks whether generating sets of fixed cardinality can be chosen so that these Cayley graphs form a superexpander family.

To describe the terminology, a family of finite graphs \((X_m)\) is an expander family if its vertex counts tend to infinity, its degrees are uniformly bounded, and there exists \(h>0\) such that
\[
|\partial_E A|\geq h|A|
\quad\text{whenever }A\subseteq V(X_m),\quad
0<|A|\leq |V(X_m)|/2,
\]
uniformly in \(m\), where \(\partial_E A\) consists of edges with exactly one endpoint in \(A\). The usual coarse-geometric superexpander requirement adds failure of uniform coarse embedding into every uniformly convex Banach space. Here a Banach space \(B\) is uniformly convex if, for every \(0<\varepsilon\leq2\), there exists \(\delta>0\) such that
\[
\|u\|,\|v\|\leq1,\quad\|u-v\|\geq\varepsilon
\quad\Longrightarrow\quad
\left\|\frac{u+v}{2}\right\|\leq1-\delta.
\]
A uniform coarse embedding of \((X_m)\) into \(B\) means maps \(f_m:V(X_m)\to B\) satisfying the coarse-embedding inequalities with the same control functions \(\rho_-,\rho_+\) for every \(m\).

The supplied source does not state its precise superexpander convention; in particular, a uniform nonlinear Poincar\'e-inequality requirement must not be silently substituted for the coarse nonembedding requirement. Nor does it specify whether \((n_m,p_m)\) is prescribed or chosen, whether the assertion is universal over such sequences, which parameters grow, or explicitly require \(|G_m|\to\infty\). Symmetry of the generating sets is likewise a convention used here to explain undirected Cayley graphs, not an additional verified source hypothesis.
\end{enumerate}

All four items are statement-only drafts. The supplied record labels the compound problem partially solved, but it does not supply precise statements of the asserted affirmative subcases.
\subsection{Short English statement}
The record collects four questions: whether \(L_p([0,1])\) coarsely embeds into \(\ell_p\); whether a free measure-preserving free-group action on a manifold has a warped cone with suitable walls; whether a slowly growing box space can have bounded successive subgroup indices; and whether special linear groups over prime fields admit generating sets of fixed cardinality producing superexpanders. Missing exponent restrictions, growth conditions, geometric conventions, and sequence quantifiers prevent a single fully specified formulation. The supplied partially solved label and literature assessment require independent review.
\subsection{Sources}
\begin{itemize}
\item Supplied local UnsolvedMath registry record, permanent ID \(30003788\), title ``Collected Coarse-Geometric Embedding Problems,'' source code \texttt{catalogue entry}, in the Oberwolfach Reports Open Problems collection. The record gives proposed year \(2018\) and names no individual proposer.
\item Supplied report citation: \emph{Mini-Workshop: Superexpanders and Their Coarse Geometry (2019)}, Oberwolfach Reports, pp.\ 1117--1160, DOI \texttt{10.4171/owr/2018/19}. Supplied source link: \url{https://doi.org/10.4171/owr/2018/19}. The citation's \(2019\) year and the record's \(2018\) proposed year are retained without resolving their relationship.
\item Supplied literature source, ``workshop report'': \url{https://ems.press/content/serial-article-files/46739}.
\item Supplied literature source, ``straightening warped cones'': \url{https://arxiv.org/abs/1705.06725}.
\item Supplied literature source, ``coarse embeddings into \(\ell_q\)'': \url{https://arxiv.org/abs/1210.0588}.
\item Literature assessment dated exactly \texttt{2026-08-21}: this is supplied dated curation, not independently verified. It labels the compound problem partially solved, reports affirmative subcases involving selected prime-power quotients and profinite-completion warped cones, and describes the manifold wall-structure, slowly-growing-box-space, general prime-field superexpander, and \(L_p\)-to-\(\ell_p\) questions as remaining open. The supplied background explicitly calls this a literature triage and requests independent verification of time-sensitive status and cited claims. It does not specify the hypotheses or conclusions of the reported affirmative subcases.
\item The registry labels its statement \texttt{corrected\_verified} and states that mathematical content was preserved while notation, typography, and omitted definitions were repaired. This is a supplied registry qualification, not independent mathematical or bibliographic verification of this draft. No external source was visited in preparing it.
\end{itemize}
\end{document}
